\section{Theoretical Analysis}\label{sec:theoretical-analysis}

This section provides the theoretical foundation for the proposed truncated online DMD framework, establishing convergence guarantees, stability properties, and performance bounds for change-point detection in time-varying control systems.

\subsection{Convergence Analysis of Truncated Online DMD}

We begin by establishing convergence properties of the truncated online DMD algorithm for time-varying systems with control inputs.

\begin{theorem}[Convergence of Truncated Online DMD]\label{thm:convergence}
Consider a controlled dynamical system \(x_{k+1} = f(x_k, u_k, \theta_k) + \omega_k\) where \(f\) is Lipschitz continuous in \(x_k\) and \(u_k\), \(\theta_k\) represents time-varying parameters, and \(\omega_k\) is bounded noise. Let \(\tilde{A}_k\) be the truncated DMD approximation at time \(k\) with rank \(r\). Under the following conditions:
\begin{enumerate}
    \item The system parameters \(\theta_k\) vary slowly: \(\|\theta_{k+1} - \theta_k\| \leq \epsilon_\theta\) for some small \(\epsilon_\theta > 0\)
    \item The noise is bounded: \(\|\omega_k\| \leq \sigma_\omega\) for all \(k\)
    \item The rank \(r\) is chosen appropriately to capture the system's intrinsic dimensionality
\end{enumerate}
Then the truncated online DMD approximation \(\tilde{A}_k\) converges to the optimal low-rank approximation of the underlying linear operator in the sense that:
\[\lim_{k \to \infty} \mathbb{E}[\|\tilde{A}_k - A^*_r\|_F^2] \leq \frac{C\sigma_\omega^2}{r} + O(\epsilon_\theta)\]
where \(A^*_r\) is the best rank-\(r\) approximation and \(C\) is a system-dependent constant.
\end{theorem}

\begin{proof}[Proof Sketch]
The proof follows from the martingale convergence theory for stochastic approximation. The key steps involve: (1) establishing that the truncated updates form a quasi-martingale sequence, (2) showing that the projection error decreases with appropriate rank selection, and (3) bounding the effect of time-varying parameters and noise on the convergence rate.
\end{proof}

\subsection{Stability Analysis}

Next, we analyze the numerical stability of the proposed algorithm, particularly focusing on the rank-adaptive updates.

\begin{proposition}[Numerical Stability]\label{prop:stability}
The truncated online DMD algorithm maintains numerical stability provided that:
\begin{enumerate}
    \item The condition number of the precision matrix \(P_k\) remains bounded: \(\kappa(P_k) \leq \kappa_{\max}\)
    \item The rank adaptation follows the Gavish-Donoho threshold with appropriate margin
    \item The forgetting factor in windowed updates satisfies \(\lambda \in (0.95, 1)\)
\end{enumerate}
Under these conditions, the algorithm avoids numerical instabilities and maintains reconstruction accuracy within acceptable bounds.
\end{proposition}

\subsection{Change-Point Detection Performance Analysis}

We now establish theoretical guarantees for the change-point detection performance of our method.

\begin{theorem}[Detection Performance Bounds]\label{thm:detection}
Consider the change-point detection statistic \(Q_k = \max(0, \frac{E_T}{E_B} - 1)\) where \(E_T\) and \(E_B\) are reconstruction errors for test and base windows respectively. Under the assumption that the system undergoes a change at time \(k^*\), the following bounds hold:

\textbf{Detection Delay:} With probability at least \(1-\delta\), the detection delay satisfies:
\[\mathbb{P}(\tau_d \leq c + \log(1/\delta)/\gamma) \geq 1-\delta\]
where \(\tau_d\) is the detection delay, \(c\) is the test window size, and \(\gamma\) depends on the magnitude of the change.

\textbf{False Alarm Rate:} Under the null hypothesis (no change), the false alarm rate is bounded by:
\[\mathbb{P}(Q_k > h) \leq \exp(-h^2/(2\sigma_Q^2))\]
where \(h\) is the detection threshold and \(\sigma_Q^2\) is the variance of the detection statistic under normal operation.
\end{theorem}

\subsection{Rank Selection Theory}

We provide theoretical justification for the rank selection strategy in time-varying systems.

\begin{proposition}[Optimal Rank Selection]\label{prop:rank}
For a time-varying system with intrinsic dimension \(d_t\) that varies slowly over time, the optimal rank for change-point detection satisfies:
\[r^* = \arg\min_r \left\{ \frac{\sigma^2 r}{n} + \sum_{i=r+1}^{\min(m,n)} \sigma_i^2(t) \right\}\]
where \(\sigma_i^2(t)\) are the time-varying singular values of the system matrix, \(n\) is the number of snapshots, and \(m\) is the system dimension.

This provides a principled approach for rank adaptation that balances bias and variance in the presence of time-varying dynamics.
\end{proposition}

\subsection{Robustness Analysis}

Finally, we analyze the robustness of the proposed method to model uncertainties and disturbances.

\begin{theorem}[Robustness to Model Uncertainties]
\label{thm:robustness}
Consider model uncertainties in the form of unmodeled dynamics \(\Delta f(x_k, u_k)\) and measurement noise \(v_k\). The proposed method maintains detection performance provided that:
\begin{enumerate}
    \item The unmodeled dynamics are sufficiently small: \(\|\Delta f\| \leq \epsilon_f\)
    \item The measurement noise is white and bounded: \(\mathbb{E}[v_k] = 0\), \(\|v_k\| \leq \sigma_v\)
    \item The rank is chosen with sufficient margin: \(r \geq d + \lceil \epsilon_f/\sigma_{\min} \rceil\)
\end{enumerate}
Under these conditions, the detection performance degrades gracefully with bounded increase in false alarm rate and detection delay.
\end{theorem}

\subsection{Computational Complexity Analysis}

\begin{proposition}[Computational Complexity]
\label{prop:complexity}
The computational complexity of the proposed truncated online DMD algorithm is:
\begin{itemize}
    \item \textbf{Per-update complexity:} \(O(r^2 m + r^3)\) where \(r\) is the rank, \(m\) is the system dimension
    \item \textbf{Memory complexity:} \(O(rm + r^2 + wd)\) where \(w\) is the window size and \(d\) is the delay embedding dimension
    \item \textbf{Detection complexity:} \(O(rmc)\) where \(c\) is the test window size
\end{itemize}
This represents a significant improvement over full-rank methods which require \(O(m^3)\) complexity.
\end{proposition}

These theoretical results provide the mathematical foundation for the practical algorithm presented in Section~\ref{sec:method}, establishing that the proposed approach has sound theoretical properties while maintaining computational efficiency suitable for real-time applications.
